% AST2000 Part 3: the habitable zone, choosing a destination, and launching at any time.
% Compile: cd "Project 3/paper" && latexmk -pdf part3.tex   (or: python build_papers.py 3 from the repo root)
% Every number written as \val{key} comes from results.tex, made by pro_3.py.
\documentclass[aps,prd,twocolumn,nofootinbib,floatfix,superscriptaddress]{revtex4-2}
\usepackage{../../common/paper/ast2000}

\begin{document}

\title{Part 3: The Habitable Zone, Choosing a Destination, and Launching at Any Time}
\author{Sakarias Frette}
\affiliation{AST2000 Rerun, University of Oslo}
\date{\today}

\begin{abstract}
We prepare the journey.
Treating the star (\val{star-T}\,K, \val{star-Rsun}\,$R_\odot$) and the planets as black bodies, we find that liquid water is possible from \val{hz-in} to \val{hz-out}\,AU.
Planets \val{hz-planets} lie inside this habitable zone.
We compare every planet as a destination by temperature, gravity, atmosphere, fuel cost and travel time, and choose planet \val{dest}.
It is rocky, has an equilibrium temperature of \val{dest-T}\,K, and costs \val{dest-dv}\,km/s and \val{dest-days} days by a Hohmann transfer, with launch windows every \val{dest-synodic}\,yr.
A $\Delta v$ map shows that burning deep in the planets' gravity wells cuts the cost to about \val{dest-oberth-dv}\,km/s.
We find how close counts as ``arrived'' (\val{dest-capture}\,km) and how close the planet must be to fill a camera pixel.
Finally we generalise the Part 1 launch to any time and any site on the equator.
All \val{n-tests} test launches were verified by \code{ast2000tools}, at most \val{max-miss}\,km from the API's own result.
\end{abstract}

\maketitle

% =====================================================================
\section{Introduction}
Part 3 of the project~\cite{ast2000project} answers three questions before the journey.
\begin{itemize}
\item Where could a lander live (\cref{sec:hz})?
\item Which planet should we go to, and when (\cref{sec:dest})?
\item Can we launch from anywhere at any time, accurately enough to plan with (\cref{sec:launch})?
\end{itemize}
Solar system seed: \val{seed}.

% =====================================================================
\section{The habitable zone}\label{sec:hz}
\subsection{Theory}
A star of radius $R_\star$ and temperature $T_\star$ radiates $L=4\pi R_\star^2\sigma T_\star^4$.
At distance $r$ this gives the flux $F=\sigma T_\star^4(R_\star/r)^2$.
A planet of radius $R$ catches $F\pi R^2$ on its cross-section and radiates $4\pi R^2\sigma T^4$ from its whole surface~\cite{carroll2017}.
In equilibrium
\begin{equation}
  T = T_\star\sqrt{\frac{R_\star}{2r}}, \qquad r(T)=\frac{R_\star}{2}\left(\frac{T_\star}{T}\right)^2. \label{eq:T}
\end{equation}
The planet's own size cancels.
Requiring liquid water, $\val{T-water-lo}\le T\le\val{T-water-hi}$\,K, defines the habitable zone through $r(T)$.
A lander needing \val{panel-power}\,W from panels of \val{panel-eff}\% efficiency needs an area $A=P/(\eta F)$.

\subsection{Results}
The star has $L=\val{star-L}\,L_\odot$.
The habitable zone spans \val{hz-in}--\val{hz-out}\,AU, or \val{hz-wide-in}--\val{hz-wide-out}\,AU with the $\pm\val{T-margin}$\,K uncertainty (\cref{fig:T,fig:hzmap}).
The home planet sits at \val{T-home}\,K and planet \val{dest} at \val{dest-T}\,K.

Planet \val{dest}'s orbit is nearly circular, so over its year the flux varies only between \val{dest-F-min} and \val{dest-F-max}\,W/m$^2$.
The temperature varies only between \val{dest-T-min} and \val{dest-T-max}\,K (\cref{fig:power}).
The lander needs a panel of just \val{dest-panel}\,m$^2$.

\widefig{1a_temperatures}{Black-body temperature against distance, with every planet (error bars: perihelion to aphelion) and the habitable zone.}{fig:T}
\widefig{1b_habitable_zone_map}{The habitable zone as a ring in the system.}{fig:hzmap}
\widefig{1c_solar_power}{Left: solar-panel area for \val{panel-power}\,W against distance. Right: planet \val{dest}'s flux over one orbit, with the matching equilibrium temperature on the right axis.}{fig:power}

% =====================================================================
\section{Choosing the destination}\label{sec:dest}
\subsection{Method}
For each planet we compute
\begin{itemize}
\item temperature;
\item surface gravity $g$ and escape speed;
\item atmospheric density;
\item day length;
\item solar-panel area;
\item the cost of getting there.
\end{itemize}
For the cost we use a Hohmann transfer: the cheapest two-burn transfer between circular orbits $r_1$ and $r_2$~\cite{curtis2014}.
It needs
\begin{align}
  \Delta v_1 &= \sqrt{\frac{GM}{r_1}}\left(\sqrt{\frac{2r_2}{r_1+r_2}}-1\right),\\
  \Delta v_2 &= \sqrt{\frac{GM}{r_2}}\left(1-\sqrt{\frac{2r_1}{r_1+r_2}}\right),
\end{align}
and takes half the transfer ellipse's period.
The target must lead home by $\pi-n_2t_{\rm transfer}$ at launch, where $n_2$ is its mean motion.
This happens once per synodic period, $1/|1/P_1-1/P_2|$.
We find the actual windows from the exact Kepler positions.

\subsection{Results}
\Cref{fig:compare} compares the candidates.
The gas giants have huge gravity and no surface to land on.
Planet 7 is hot and close to the star, and it costs a lot to reach.
Planet \val{dest} is the best compromise.
\begin{itemize}
\item It is rocky and inside the habitable zone.
\item Its gravity is $g=\val{dest-g}$\,m/s$^2$, against \val{home-g3} at home.
\item A day lasts \val{dest-day}\,h.
\item A Hohmann transfer costs $\val{dest-dv1}+\val{dest-dv2}=\val{dest-dv}$\,km/s and takes \val{dest-days} days.
\end{itemize}
Our automatic rule (rocky, habitable if possible, then cheapest) picks planet \val{dest-auto} as well.
Its atmosphere is dense, \val{dest-rho-atm}\,kg/m$^3$ at the surface.
That is a challenge for the landing, but also a free brake.

The window calendar (\cref{fig:windows}) gives windows at $t=\val{dest-windows},\ldots$\,yr.
Part 5 uses the second.

The $\Delta v$ map (\cref{fig:dvmap}) puts the choice in context.
Going outward, costs rise quickly at first and then level off.
The real trip can be much cheaper than the orbit-to-orbit Hohmann estimate.
Right after the Part 1 launch the rocket is still deep in the home planet's gravity well, moving at the local escape speed.
There, a burn adds much more energy per unit of $\Delta v$ (the Oberth effect), so leaving with $v_\infty$ costs only $\sqrt{v_\infty^2+v_{\rm esc}^2}-v_{\rm esc}$.
At the other end, capture into the cheapest circular orbit costs $v_\infty/\sqrt2$.
Together these give about \val{dest-oberth-dv}\,km/s for planet \val{dest}, which Part 5 confirms in the full simulation.

\subsection{How close is close?}
The spacecraft counts as arrived when the planet pulls $k=\val{k-dominant}$ times harder than the star, which happens inside
\begin{equation}
  l = |\vb r|\sqrt{\frac{M_p}{k M_\star}}.
\end{equation}
For planet \val{dest} this is \val{dest-capture}\,km, \val{dest-capture-radii} planet radii (\cref{fig:close}).
A camera with $P=\val{camera-px}$ pixels across a $\val{camera-fov}^\circ$ field of view resolves the planet once $2R/L>F/P$.
That holds within \val{dest-resolve}\,km, so we see the planet long before its gravity dominates.

\widefig{2a_destination_comparison}{Every candidate compared. Hatched: gas giants. HZ: inside the habitable zone.}{fig:compare}
\colfig{2b_transfer_routes}{Hohmann routes from home to the rocky planets.}{fig:routes}
\widefig{2c_window_calendar}{Top: the angle by which planet \val{dest} leads home, against the angle a Hohmann transfer needs. A window opens where they meet. Bottom: windows to every rocky planet.}{fig:windows}
\widefig{2d_delta_v_map}{$\Delta v$ map. Hohmann cost from orbit to orbit (black), and the cost with an Oberth departure right after launch plus the cheapest capture (violet). Right: travel time.}{fig:dvmap}
\widefig{4_capture_and_camera}{Capture distance ($k=\val{k-dominant}$) and camera-resolution distance for every planet.}{fig:close}

% =====================================================================
\section{Launching at any time}\label{sec:launch}
\subsection{Method}
Seen from the planet, the launch is the same at any time and from any site: radially up, \val{drift-km}\,km of sideways drift from the rotation.
Only the step to the star's frame changes.
With $\hat{\vb e}_r=(\cos\phi,\sin\phi)$ at site angle $\phi$ and $\hat{\vb e}_t$ along the rotation,
\begin{align}
  \vb r &= \vb r_p(t_0+\Delta t) + v_{rot}\Delta t\,\hat{\vb e}_t + r_{\rm final}\hat{\vb e}_r,\\
  \vb v &= \vb v_p(t_0+\Delta t) + v_{rot}\hat{\vb e}_t + v_r\hat{\vb e}_r.
\end{align}
The planet's state comes from Kepler's equation (Part 2), not from interpolating our simulated orbits.
Those drift hundreds of kilometres from the truth over 20 home orbits, while the API's tolerance is only \val{tolerance}\,km (1.5\% of the planet radius).
Using the curved orbit instead of Part 1's straight line changes the result by \val{curved-vs-straight}\,km.

\subsection{Results}
We tested \val{n-tests} launches spread over 20 home orbits and all directions (\cref{fig:sites,fig:times}).
\val{n-verified} of \val{n-tests} were verified.
The largest miss was \val{max-miss}\,km in position and \val{max-miss-v}\,m/s in velocity.
The velocity difference comes from the API's 0.01\,s Euler steps, and Part 5 corrects for it after launch.

\widefig{3a_launch_sites}{Launches from eight sites on the equator (planet frame), and a zoom on the sideways drift from the rotation.}{fig:sites}
\widefig{3b_launch_times}{Test launches along the home orbit, and the distance between our predicted position and the API's.}{fig:times}

% =====================================================================
\section{Discussion and conclusion}
The black-body model ignores albedo and greenhouse heating, both of which would move the zone.
The $\pm\val{T-margin}$\,K band shows how sensitive the edges are.
Our choice is robust: planet \val{dest} sits well inside the zone either way.

The Hohmann estimate is deliberately simple: circular, coplanar orbits and no planetary gravity.
The $\Delta v$ map shows that it overestimates the real cost by more than a factor of two, because it ignores the Oberth effect.
That is the most important lesson carried into the journey.

The launch is now a tool we can use at any time and site with kilometre accuracy.
That is exactly what the journey planner in Part 5 needs.

\bibliography{../../common/paper/references}

\appendix
\section{Mission log}
\begin{missionbox}{Part 3 key numbers (seed \val{seed})}
Habitable zone & \val{hz-in}--\val{hz-out}\,AU \\
Destination & planet \val{dest} (\val{dest-type}, $a=\val{dest-a}$\,AU) \\
Temperature / panel & \val{dest-T}\,K / \val{dest-panel}\,m$^2$ \\
Hohmann $\Delta v$ / time & \val{dest-dv}\,km/s / \val{dest-days}\,d \\
With Oberth & \val{dest-oberth-dv}\,km/s \\
Windows (yr) & \val{dest-windows} \\
Capture distance & \val{dest-capture}\,km \\
Launch tests verified & \val{n-verified}/\val{n-tests} \\
\end{missionbox}

\end{document}
